\section{Discussion}

\subsection{Stress as Task-Dependent Behavioural Modulation}
Our goal is not to present a deployed stress-adaptive VR system but to characterise how acute stress changes behaviour across interaction components. The MRT-informed task-probe framing offers an analytical basis for reasoning about stressful VR situations that have not yet been directly studied. By considering the resource-demand profile of the current task, designers can anticipate which aspects of interaction may remain stable and which may require support. Our findings suggest that acute stress in VR should not be treated as a uniform disruptor of interaction in VR. Across the three task probes, stress-related changes appeared selectively across measures rather than as a consistent decline in speed, accuracy, or efficiency.

For the target acquisition task, we did not detect statistically significant phase differences in movement time, throughput, or endpoint offset. These findings do not establish that target acquisition is unaffected by stress. The exploratory pre-hit analysis describes a possible phase-associated change in terminal offset variability that requires further investigation. 
%Exploratory axis-wise analyses further indicated that stress-related changes in endpoint error may be direction-specific rather than isotropic: the vertical component showed a modest increase from \textit{Baseline} measurements to \textit{Post-Stress} although the analysis of overall offset did not detect a statistically significant phase effect. We interpret this cautiously, but it raises the possibility that VR pointing under stress may be modulated by axis-specific factors such as hand-tracking characteristics, posture, or embodied spatial alignment.  

For the reading aloud task, stress effects were more evident in vocal characteristics than in the task itself. Mean F0 increased after stress induction, whereas no statistically significant phase effects were detected in reading duration, speech rate, and active speech ratio remained comparatively stable. Constrained verbal production may therefore preserve coarse temporal structure while still reflecting stress through prosodic arousal.

For the ToL task, stress influenced the total duration and execution time, but did not have no statistics significance on first-move latency and move counts. We do not interpret this as evidence that stress improved planning ability. A more cautious reading is that stress altered execution dynamics after the initial planning period, possibly reflecting a more action-oriented or urgent interaction style. This phenomenon may be consistent with attentional narrowing accounts of stress and arousal, suggesting that stress can bias attention toward immediate task-relevant goals while reducing the breadth of monitoring or deliberation~\cite{easterbrook1959effect}. 
Thus, participants may have completed the ToL task more quickly under stress not because planning improved, but because they shifted toward more immediate, action-oriented execution.

These results indicate that acute stress in immersive environments may manifest less as a consistent global decrement and more as selective behavioural changes across tasks with different resource-demand profiles. This interpretation helps explain why some metrics remained stable while others showed more localised or phase-specific effects. One possible explanation comes from Multiple Resource Theory (MRT)~\cite{wickens2008multiple}, which suggests that performance depends on multiple partially independent resource channels spanning perceptual, cognitive, and response processes. From this perspective, acute stress need not reduce all forms of performance equally: it may instead manifest differently depending on the resource demands of the task. Our results are broadly consistent with such an account. Tasks relying more strongly on higher-level planning showed clearer changes in execution dynamics, perceptual-motor performance remained relatively stable at the level of predefined outcome metrics, and verbal changes were expressed primarily through vocal arousal rather than broad timing disruption. This pattern suggests that stress reshapes interaction differently depending on the resources most relevant to a given task.
\subsection{Practical significance}

\subsection{Implications for Stress-Aware VR Design}
Our findings have several implications for the design of stress-aware immersive systems. Most importantly, stress-aware adaptation in VR should not assume that stress affects all interaction in the same way. Hence, adaptive supports may be more effective when informed by the demands of the current task.

First, stress-aware systems should adapt support to the current task rather than applying a generic stress response. Planning-heavy tasks may benefit from subgoal cues, step previews, or reduced concurrent demands, whereas simple target selection may require little intervention unless fine precision is critical. Second, systems should monitor changes in interaction dynamics, not only overt errors or failures: in our study, several stress-related changes appeared in timing, vocal arousal, or terminal movement dynamics rather than in accuracy. Third, support may need to continue into recovery, since some measures did not immediately return to baseline, suggesting that post-stress interaction may remain altered even after the stressor has ended.

More broadly, these findings point toward opportunities for stress-aware VR systems that integrate user state with task structure. We do not claim to have validated a specific adaptive system. Here, our results suggest that future adaptive immersive systems may benefit from considering not only whether a user is stressed, but also what kind of interaction they are currently performing.

\subsection{Limitations and Future Work}
A first limitation concerns potential learning and order effects.Because phase order was fixed, phase-associated changes cannot be attributed specifically to stress, as practice, fatigue, and familiarisation may also have contributed. Although we used a Latin-square-based ordering strategy and varied the reading aloud passages and ToL configurations across rounds, the repeated-measures design may still have introduced practice-related changes. This issue was considered during pilot testing, where repeated exposure to multiple ToL and target acquisition trials within the same session appeared to amplify familiarity effects, informing our decision to use a more constrained design in the main study with one ToL instance and one target acquisition block per round. Even so, residual practice effects cannot be fully ruled out. Our sample size is comparable to many controlled within-subject HCI studies, but may still limit sensitivity to modest task-specific effects and individual-difference patterns.

A second limitation is that our three tasks were intentionally designed as controlled probes rather than as full end-to-end VR workflows. Although target acquisition, reading aloud, and ToL task allowed us to isolate specific classes of interaction under stress, they do not fully capture the complexity, context switching, or social and semantic richness of real-world VR use. Our findings should therefore be interpreted as evidence about stress-sensitive interaction components rather than as a complete account of stress effects in deployed VR applications.

A third limitation concerns the scope of our MRT-based interpretation. Although the task battery was informed by MRT, the study did not directly test resource competition in the classical MRT sense: we did not employ dual-task or concurrent-task paradigms, nor did we manipulate competing demands within the same trial. Instead, we compared separate task probes that differed in their dominant interaction demands. MRT should therefore be understood here as a task-selection and interpretive lens, rather than as a theory directly validated by our experimental design.

A final limitation is that TSST induction is unlikely to operate as a uniform stress dose across individuals. Prior work suggests that responses to the TSST may depend on participant-specific factors such as ability, engagement, and coping style~\cite{kudielka2007ten}, which is especially relevant in the arithmetic phase, where participants may differ in response speed, accuracy, and how they react to errors or interruptions. One advantage of VR is that such induction-phase behaviour can be both tightly controlled and continuously logged. We explored whether arithmetic-phase dynamics were associated with downstream task changes, but did not observe robust significant relationships in the current sample. Future work should examine these individual-difference effects more systematically with larger samples. Our participants were young adults recruited from a university population, so the findings may not generalise to older users or to individuals with cognitive impairment, elevated baseline anxiety, anxiety-related conditions, or other psychosocial factors that could influence stress reactivity and interaction performance.



